\section{Lambert inversion and exact logarithmic coefficients}
\label{lmi:section}

The marginal law at $s=3$ in Theorem~\ref{qfb:phase} converges slowly.
Its logarithmic corrections nevertheless have exact finite formulas.
We first reconcile the order-three inverse chart and correction calculus
from \cite{qfb:reports}. A Laplace representation then combines the
correction calculus into a single Gamma moment at each order. The same
representation gives an exact convolution identity, while the chart
coefficients have a Catalan principal part.

\subsection{An exact marginal chart and its rational coefficients}
Throughout this section put
\[
 A=\zeta(2),\quad u=U_c(3),\quad v_c=u-A,\quad R=e^{-v_c},
 \quad d=-\operatorname{Log}(r/R),\quad x=-\operatorname{Log}(re^v).
\]
Thus $v=v_c+d-x$, and \eqref{qfb:normal} becomes
\begin{equation}\label{lmi:normal}
 (A+x)d-\frac{d^2}{2}
 =\frac{x^2}{2}\left(\log\frac1x+\frac52\right)+x^3C_3(x),
 \qquad
 C_3(x)=\sum_{j\ge0}\frac{(-1)^{j+3}\zeta(-j)}{(j+3)!}x^j.
\end{equation}
The tail is holomorphic for $|x|<2\pi$, with
$C_3(x)=1/12-x/288+O(x^3)$. All branches below continue from
small positive $d$. Let $W_{-1}$ denote the lower real Lambert
branch \cite{lmi:lambert}, and set
\begin{equation}\label{lmi:core}
 \Lambda=-\frac12W_{-1}(-4Ad e^{-5}),\qquad
 \tau=\sqrt{\frac{2Ad}{\Lambda}}.
\end{equation}
Exactly, $\Lambda=\log(1/\tau)+5/2$ and $\tau^2\Lambda=2Ad$.

\begin{theorem}[Convergent Lambert chart]\label{lmi:chart}
There is a holomorphic germ $U(p,w)$, $U(0,0)=1$, determined by
\begin{equation}\label{lmi:chart-equation}
 U^2(1-w\log U)+2pw^2U^3C_3(pwU)
 =1+\frac{pw}{A}U-\frac{p^2w}{4A^2}.
\end{equation}
The selected inverse of \eqref{lmi:normal} is exactly
$x=\tau U(\tau\Lambda,\Lambda^{-1})$. Moreover, $U-1$ is divisible
by $pw$ and
\begin{equation}\label{lmi:first-correction}
 x=\tau+\tau^2\frac{\Lambda/A-1/6}{2\Lambda-1}
       +O(\tau^3\Lambda).
\end{equation}
\end{theorem}
\begin{proof}
Substitute $x=\tau U$, $p=\tau\Lambda$, $w=\Lambda^{-1}$,
$\tau=pw$, and $d=p^2w/(2A)$ in \eqref{lmi:normal}.
Division by $Ad$ gives \eqref{lmi:chart-equation}.
The derivative in $U$ at $(0,0,1)$ is $2$, so the holomorphic
implicit function theorem applies. At either $p=0$ or $w=0$,
uniqueness gives $U=1$. The coefficient of $p$ is
\[
 [p]U=\frac{w/A-w^2/6}{2-w};
\]
the remaining terms are $O(p^2w)$. This proves
\eqref{lmi:first-correction}. The positive solution and then analytic
continuation identify the inverse already selected in Theorem~\ref{qfb:phase}.
\end{proof}

Writing $U=1+\sum_{n\ge1}u_n(w)p^n$, define
\[
 F(h,w)=(1+h)^2\{1-w\log(1+h)\}-1
\]
and
\[
 \Psi(p,w,h)=\frac{pw}{A}(1+h)-\frac{p^2w}{4A^2}
             -2pw^2(1+h)^3C_3(pw(1+h)).
\]
An auxiliary multiplier on $\Psi$ and Lagrange inversion give the finite
identity
\begin{equation}\label{lmi:finite}
 u_n(w)=\sum_{k=1}^n\frac1k[p^nh^{k-1}]
                 \left\{\frac{h\Psi(p,w,h)}{F(h,w)}\right\}^k.
\end{equation}
Only $k\le n$ contributes because $\Psi$ is divisible by $p$.

\begin{proposition}[Exact Catalan principal parts]\label{lmi:poles}
Each $u_n(w)$ is rational in $w$, with coefficients in $\mathbb Q(A)$.
Its only finite pole is at $w=2$, has exact order $2n-1$, and satisfies
\begin{equation}\label{lmi:catalan}
 \lim_{w\to2}(2-w)^{2n-1}u_n(w)
 =2^{n-1}\operatorname{Cat}_{n-1}
       \left(\frac2A-\frac23\right)^n .
\end{equation}
\end{proposition}
\begin{proof}
In \eqref{lmi:finite}, the only denominators involving $w$ arise from
$F/h=(2-w)+(1-3w/2)h+O(h^2)$. In the summand with index $k$,
extracting $h^{k-1}$ produces at most $2k-1$ factors of $(2-w)^{-1}$.
The largest pole can therefore come only from $k=n$, taking $p$
once from every $\Psi$ and every $h$ from the denominator expansion.
At $w=2$, the linear coefficient of $F/h$ is $-2$ and
$[p]\Psi(p,2,0)=2/A-2/3$. The resulting factor is
$2^{n-1}\binom{2n-2}{n-1}/n$, which proves \eqref{lmi:catalan}.
It is nonzero since $A=\zeta(2)<2$. This is a statement about the
rational coefficient functions; the chart is proved near $w=0$ and
no global physical singularity at $w=2$ is inferred.
\end{proof}

\subsection{A Laplace identity and an exact convolution}
The Lambert core has a useful identity beyond its asymptotic expansion.
For real $H>1$, let $T(H)>1$ be the solution of
$T-\log T=H$, equivalently $T=-W_{-1}(-e^{-H})$.
For $\nu>0$ and $t>0$ define
\begin{equation}\label{lmi:density}
 \rho_\nu(t)=\frac{\nu\,t^{t+\nu-1}}{\Gamma(t+\nu+1)}.
\end{equation}
Connections between the lower Lambert branch, Laplace transforms and
probability laws are classical; see \cite{lmi:pakes}. We give a direct
proof in the normalization needed here, without a priority claim.

\begin{theorem}[Gamma density identities]\label{lmi:laplace}
For $\nu>0$ and $H>1$,
\begin{equation}\label{lmi:laplace-identity}
 T(H)^{-\nu}=\int_0^\infty e^{-Ht}\rho_\nu(t)\,dt .
\end{equation}
For every $\mu,\nu,t>0$,
\begin{equation}\label{lmi:convolution}
 \int_0^t\rho_\mu(y)\rho_\nu(t-y)\,dy=\rho_{\mu+\nu}(t).
\end{equation}
In particular, $e^{-t}\rho_\nu(t)$ integrates to one, and these
densities form a convolution family in $\nu$.
\end{theorem}
\begin{proof}
The map $z\mapsto z-\operatorname{Log}z$ is one-to-one on $\Re z>1$:
its derivative has positive real part, and integrating that derivative
along a line segment proves injectivity. Its boundary has real part
$1-\tfrac12\log(1+y^2)\le1$ and imaginary part $y-\arctan y$,
strictly increasing onto the real line. Also its modulus tends to
infinity with $|z|$. The argument principle, or continuation from the
positive axis in this boundary description, therefore gives a
holomorphic inverse $T(H)$ for $\Re H>1$, with
$T(H)=H+\operatorname{Log}H+o(1)$ as $|\Im H|\to\infty$
on each fixed vertical line.

Initially take $\nu>1$, so Bromwich inversion of $T(H)^{-\nu}$
is absolutely convergent. Set $H=z-\operatorname{Log}z$ in that
integral. The image contour has
$\Re z=c+\log|z|$, where $c>1$ is the original Bromwich abscissa.
It can be deformed to a fixed vertical line $\Re z=b>1$.
On the horizontal connectors the integrand is
$O(|\Im z|^{-\nu})$ and their lengths are $O(\log|\Im z|)$;
their contributions tend to zero. Inverse Laplace inversion for
powers now gives
\begin{align*}
 \frac1{2\pi i}\int e^{tz}z^{-t-\nu}(1-z^{-1})\,dz
 &=\frac{t^{t+\nu-1}}{\Gamma(t+\nu)}
   -\frac{t^{t+\nu}}{\Gamma(t+\nu+1)}\\
 &=\rho_\nu(t).
\end{align*}
This proves \eqref{lmi:laplace-identity} for $\nu>1$ by Laplace
uniqueness. Near zero, $\rho_\nu(t)\sim t^{\nu-1}/\Gamma(\nu)$;
Stirling's formula gives $\rho_\nu(t)\sim\nu e^t/
(\sqrt{2\pi}\,t^{3/2})$ at infinity. The integral is thus holomorphic
in $\nu$ on $\Re\nu>0$, locally uniformly, and analytic continuation
proves the assertion for every positive $\nu$.

Tonelli's theorem and \eqref{lmi:laplace-identity} show that the
Laplace transform of the convolution in \eqref{lmi:convolution}
is $T(H)^{-\mu-\nu}$. Laplace uniqueness gives equality almost
everywhere, and continuity for $t>0$ gives equality everywhere.
Finally let $H\downarrow1$ and use monotone convergence and $T(H)\downarrow1$.
\end{proof}

\subsection{Convergent core polynomials and asymptotic coefficient polynomials}
Put
\[
 L=\log(1/d),\qquad c=5-\log(4A),\qquad T=2\Lambda .
\]
Then $T-\log T=L+c$ and $\tau=2\sqrt{Ad/T}$. The report's
convergent logarithmic core expansion is
\begin{equation}\label{lmi:core-series}
 \tau=2\sqrt{Ad}\,L^{-1/2}
       \sum_{j\ge0}L^{-j}P_j(\log L+c),
\end{equation}
where $P_0=1$ and, for $j\ge1$,
\begin{equation}\label{lmi:core-polynomials}
 P_j(Y)=-\frac1{2j}[h^{j-1}]
             (1+h)^{-3/2}\{Y+\log(1+h)\}^j.
\end{equation}
Indeed $T=L(1+V)$ gives
$V=\eta+p\log(1+V)$ with $p=L^{-1}$ and $\eta=p(\log L+c)$.
The implicit germ is holomorphic in $(p,\eta)$ at zero.
Lagrange inversion after $\eta=pY$ gives
\eqref{lmi:core-polynomials}; grouping its normally convergent
double series proves \eqref{lmi:core-series}. In particular,
$P_1=-Y/2$ and $P_2=3Y^2/8-Y/2$.

For comparison with the report's finite transfer formula, set
\[
 G(a)=\frac1{\Gamma(-a)},\qquad
 M_k(a)=\frac{(-1)^kG^{(k)}(a)}{k!\,G(a)}
\]
and define
\begin{equation}\label{lmi:transfer-polynomials}
 Q_N(Y)=\sum_{j=0}^N M_{N-j}(1/2)
       \prod_{\ell=0}^{N-j-1}
       \left(\frac{d}{dY}-j-\frac12-\ell\right)P_j(Y).
\end{equation}
An empty product is one. This expression is finite.

\begin{theorem}[One Gamma moment for every logarithmic correction]
\label{lmi:moments}
The polynomials in \eqref{lmi:transfer-polynomials} have the exact formulas
\begin{align}
 Q_N(Y)
 &=\frac1{\sqrt\pi\,N!}\Re\int_0^\infty
       e^{-y}y^{N-1/2}(\log y-Y+i\pi)^N\,dy \label{lmi:moment}\\
 &=\left.\frac1{\sqrt\pi\,N!}
       \Re(\partial_z-Y+i\pi)^N\Gamma(z)\right|_{z=N+1/2}.
       \label{lmi:gamma}
\end{align}
Here $Y$ is real; the resulting real polynomial extends uniquely to
complex $Y$. Its leading coefficient is
$(-1)^N\binom{2N}{N}/4^N$, and
\begin{equation}\label{lmi:derivative-recurrence}
 Q_{N+1}'=Q_N'-(N+\tfrac12)Q_N.
\end{equation}
For $a_n=[r^n]v(r)$ in the critical order-three model and every fixed
integer $M\ge0$,
\begin{equation}\label{lmi:coefficient-law}
\begin{split}
 R^na_n={}&\sqrt{\frac A\pi}\,n^{-3/2}(\log n)^{-1/2}
 \left\{\sum_{N=0}^M
 \frac{Q_N(\log\log n+c)}{(\log n)^N}\right.\\
 &\left.\hspace{24mm}
 +O\left(\frac{(\log\log n)^{M+1}}{(\log n)^{M+1}}\right)\right\}.
\end{split}
\end{equation}
Unlike the core expansion, this coefficient expansion is asserted
asymptotically to each fixed order, not as a convergent infinite sum.
\end{theorem}
\begin{proof}
First, power--logarithm transfer \cite{qfb:transfer} applied to
\eqref{lmi:core-series} gives \eqref{lmi:transfer-polynomials}.
On a Hankel contour, $1-r/R=t/n$ changes the large logarithm to
$\log n-\log t$. Taylor expansion contributes
$(-1)^k(\log t)^k/k!$, whose integral against the power kernel
is $G^{(k)}(1/2)$. Repeated differentiation of
$L^{-j-1/2}P_j(\log L+c)$ gives exactly the product in
\eqref{lmi:transfer-polynomials}. The transfer domain was established
in Theorem~\ref{qfb:phase}. The analytic chart gives
$x-\tau=O(\tau^2)$, and $d=1-r/R+O((1-r/R)^2)$; both changes
affect only terms smaller by a positive power of $n$, up to logarithms.
Finite differentiated remainders therefore give
\eqref{lmi:coefficient-law}.

To combine the finite operations, use \eqref{lmi:laplace-identity}
with $\nu=1/2$ and $H=L+c$. It gives the exact local identity
\[
 \tau=2\sqrt A\int_0^\infty
             e^{-ct}\rho_{1/2}(t)d^{t+1/2}\,dt .
\]
For logarithmic asymptotics it suffices to retain $0\le t\le\epsilon$
with a fixed $0<\epsilon<1/4$. The tail is
$O(d^{1/2+\epsilon'})$ on each closed slit subsector, for any
$0<\epsilon'<\epsilon$, as are its appropriately scaled derivatives.
Uniform power transfer on this compact $t$-interval reduces the
leading logarithmic sector of $-\tau$ to
\[
 \frac{\sqrt A}{\pi}n^{-3/2}
 \int_0^\epsilon e^{-(\log n+c)t}t^{t-1/2}\cos(\pi t)\,dt .
\]
The simplification uses the exact reflection identity
\[
 \frac1{\Gamma(t+3/2)\Gamma(-t-1/2)}
       =-\frac{\cos(\pi t)}{\pi}.
\]
Set $t=y/\log n$. Watson expansion of
$\exp\{t(\log t-c)\}\cos(\pi t)$ gives precisely
\eqref{lmi:moment}. Only finitely many Taylor terms are used for a
fixed target order; no integral with the transferred density over the
whole positive axis is asserted. Both calculations give polynomial
asymptotic expansions of the same function, whose coefficients are
unique, proving equality with \eqref{lmi:transfer-polynomials}.

Differentiation under the convergent Gamma integral proves
\eqref{lmi:gamma}. Its highest power of $Y$ is
$(-1)^N\Gamma(N+1/2)/(\sqrt\pi\,N!)$.
Finally apply $\Gamma(z+1)=z\Gamma(z)$ in
\eqref{lmi:gamma} and differentiate with respect to $Y$.
The product rule gives \eqref{lmi:derivative-recurrence}.
\end{proof}

The centering which removes the digamma constant is especially simple.
Let $H=Y-\psi(-1/2)$ and
$\widehat Q_N(H)=Q_N(H+\psi(-1/2))$. Then
\begin{align*}
 \widehat Q_0(H)&=1,\qquad \widehat Q_1(H)=-H/2,\\
 \widehat Q_2(H)&=\frac38H^2-\frac12H-\frac{3\pi^2}{16}-\frac32,\\
 \widehat Q_3(H)&=-\frac5{16}H^3+H^2
       +\left(\frac{13}{4}+\frac{15\pi^2}{32}\right)H
       +1-\frac{\pi^2}{2}-\frac{35}{8}\zeta(3).
\end{align*}
In \eqref{lmi:coefficient-law}, the centered argument is
$H=\log\log n-\log A+\gamma+3$.
The first two corrections agree with the incoming report; the Gamma
moment formula supplies every further coefficient directly.

Independent symbolic checks compare the two formulas through $N=8$,
verify their derivative recurrence, and reproduce the exact pole orders
through $n=5$ using implicit coefficient recursion at a rational
diagnostic value of $A$. Two precisions of quadrature check the Laplace
and convolution identities, separately from those exact checks.
The conventional proofs above establish the all-order statements.
The nonmarginal inverse charts and joint crossover branches retain
separate integration scopes.
